% !TeX root = ../main.tex
\section{Companion AI}
\label{sec:ai}

\subsection{Why an AI Is Needed}

One of the barriers to an abstract board game is the need for an opponent.
A companion AI lets a player practise alone and feel the tactical depth
without having to wait for someone else.

\subsection{Why Not Random-Rollout MCTS}

Monte Carlo Tree Search (MCTS)~\citep{coulom2006efficient,browne2012survey}
is often chosen for multiplayer games because it needs no explicit
evaluation function. The argument below targets \emph{MCTS with plain
random rollouts} (vanilla UCT)---not MCTS in general: the neural-network-%
guided MCTS variant of the AlphaZero kind~\citep{silver2018alphazero}
replaces random rollouts with a learned value function and plays very
strongly even in the Reversi-like class, so it lies outside the scope of
this critique (and is revisited under future work). With plain random
rollouts and a small compute budget, MCTS performs poorly in Tactical
Formation because:
\begin{itemize}
  \item A single move can flip 5--10 pieces at once---the outcome of a
        random game is almost unrelated to the quality of the first move.
  \item The rollout signal is too noisy, and UCT cannot distinguish a good
        move from a bad one.
  \item With a budget of 500--800 simulations, MCTS explores only the first
        1--2 layers, equivalent to greedy 1-ply.
\end{itemize}
This contrasts with Hex, where MCTS rollouts are more effective because the
rollout signal is less noisy~\citep{arneson2010hex}: how well a random
rollout fits depends strongly on its reliability within each class of game.
In the Reversi-like class, an explicit evaluation function based on piece
count and mobility---in the spirit of the learned evaluation functions for
Othello~\citep{buro1999how}---distinguishes move quality better than a
random rollout.

\subsection{Adversarial Search}
\label{ssec:search}

The bot uses \emph{two} separate algorithms depending on the number of
players---both from the Minimax family with Alpha-Beta pruning, but
differing in their evaluation function and in how they handle the endgame.
We write $N_c$ for the number of pieces and $M_c$ for the number of legal
moves (mobility) of side $c$.

\paragraph{2-player mode---Minimax Alpha-Beta.} The two sides are fully
opposed, so the bot uses classic Minimax~\citep{knuth1975alphabeta}
(Algorithm~\ref{alg:minimax}): our side maximises, the opponent minimises,
alternating via the \textsc{max} flag. The evaluation function is designed
to \emph{think like a player}---both wanting to capture pieces and wanting
to squeeze the opponent's space of moves---by normalising both terms to
$[-1,1]$ and then assigning weights $0.7$ and $0.3$ (see the reasoning in
Section~\ref{ssec:pitfalls}):
\begin{equation}
\mathrm{eval}_2(b,c)=0.7\,\frac{N_c-N_{\bar c}}{N_c+N_{\bar c}}
+0.3\,\frac{M_c-M_{\bar c}}{M_c+M_{\bar c}}\ \in[-1,1].
\label{eq:eval}
\end{equation}
When a side is no longer alive, the terminal score carries a fast-win bonus
\begin{equation}
\mathrm{score}^{2}_{\text{term}}=\pm(10000+d)
\label{eq:term2}
\end{equation}
($+$ if the opponent loses, $-$ if we lose; $d$ is the remaining depth).
The terminal value dominates every eval score $\in[-1,1]$, so shrinking the
scale does not make the bot miss a real winning or losing move.

\begin{algorithm}[H]
\caption{Minimax Alpha-Beta (2-player mode)}
\label{alg:minimax}
\begin{algorithmic}[1]
\Function{Minimax}{$s,\ d,\ \alpha,\ \beta,\ \textsc{max},\ \textsc{me},\ \textsc{opp}$}
  \If{$s$ is terminal} \State \Return $\pm(10000+d)$ \Comment{by~\eqref{eq:term2}} \EndIf
  \If{$d=0$} \State \Return $\mathrm{eval}_2(s,\textsc{me})$ \EndIf
  \State $c \gets \textsc{me}$ if \textsc{max}, else $\textsc{opp}$
  \State $\mathit{moves} \gets \mathcal A(c,s)$
  \If{$\mathit{moves}=\varnothing$} \State \Return $\mathrm{eval}_2(s,\textsc{me})$ \EndIf
  \If{\textsc{max}} \Comment{our turn: maximise}
    \State $v \gets -\infty$
    \For{$a \in \mathit{moves}$}
      \State $v \gets \max\bigl(v,\ \textsc{Minimax}(T(s,a),\,d{-}1,\,\alpha,\,\beta,\,\textsf{false},\,\textsc{me},\,\textsc{opp})\bigr)$
      \State $\alpha \gets \max(\alpha, v)$
      \If{$\alpha \ge \beta$} \textbf{break} \Comment{$\beta$ cutoff} \EndIf
    \EndFor
  \Else \Comment{opponent's turn: minimise}
    \State $v \gets +\infty$
    \For{$a \in \mathit{moves}$}
      \State $v \gets \min\bigl(v,\ \textsc{Minimax}(T(s,a),\,d{-}1,\,\alpha,\,\beta,\,\textsf{true},\,\textsc{me},\,\textsc{opp})\bigr)$
      \State $\beta \gets \min(\beta, v)$
      \If{$\alpha \ge \beta$} \textbf{break} \Comment{$\alpha$ cutoff} \EndIf
    \EndFor
  \EndIf
  \State \Return $v$
\EndFunction
\end{algorithmic}
\end{algorithm}

\paragraph{3--4 player mode---Paranoid Minimax.} With $\ge 3$ sides, the
score is no longer fully opposed, so pruning is
harder~\citep{korf1991multiplayer,sturtevant2000pruning}. The bot applies
\textbf{Paranoid Minimax}~\citep{sturtevant2003comparison,sturtevant2008analysis}
(Algorithm~\ref{alg:paranoid}): the side under consideration maximises,
\emph{every other side is treated as a single coalition} that minimises,
reducing the $N$-player problem to a two-sided max--min so that full
Alpha-Beta applies. It keeps the same $0.7/0.3$ weights and the piece-count
term normalised to $[-1,1]$ as in the 2-player mode, but because the
opponents are merged into one coalition, the mobility term is made
\emph{asymmetric}: it measures only the mobility $M_c$ of the side under
consideration (without subtracting the opponents' mobility) and normalises
by the total piece count $T$ to stay within $[0,1)$:
\begin{equation}
\mathrm{eval}_{\text{par}}(b,c)=0.7\,\frac{N_c-N_{\bar c}}{T}
+0.3\,\frac{M_c}{M_c+T},\qquad T=N_c+N_{\bar c},
\label{eq:evalpar}
\end{equation}
where $N_{\bar c}$ is the total piece count of \emph{all} opponents. Unlike
\eqref{eq:eval}, the mobility term here is never negative and is never
compared directly with an opponent; it only rewards a side that still has
many moves (avoiding a blockade)---a pragmatic choice, since the mobility
of a merged coalition does not clearly reflect the threat from each
individual opponent. When only one side remains alive, the terminal score
carries a fast-win bonus
\begin{equation}
\mathrm{score}^{\text{par}}_{\text{term}}=\pm(10000+d),
\label{eq:termpar}
\end{equation}
where $d$ is the remaining depth---encouraging the bot to finish early
rather than prolong the game. A strong alternative is Best Reply
Search~\citep{schadd2011brs}, which often outperforms Paranoid in many
4-player games; we choose Paranoid for its simplicity and because it is good
enough within a real-time compute budget, leaving a comparison with Best
Reply Search to future work.

\begin{algorithm}[H]
\caption{Paranoid Minimax (3--4 player mode)}
\label{alg:paranoid}
\begin{algorithmic}[1]
\Function{Paranoid}{$s,\ d,\ \alpha,\ \beta,\ \textsc{me}$}
  \If{only one side is alive} \State \Return $\pm(10000+d)$ \Comment{by~\eqref{eq:termpar}} \EndIf
  \If{$d=0$} \State \Return $\mathrm{eval}_{\text{par}}(s,\textsc{me})$ \EndIf
  \State $c \gets \textsc{turn}(s)$;\quad $\mathit{moves} \gets \mathcal A(c,s)$
  \If{$\mathit{moves}=\varnothing$} \State \Return \textsc{Paranoid}(drop $c$, pass turn) \Comment{a blockaded side is skipped} \EndIf
  \If{$c = \textsc{me}$} \Comment{maximising layer}
    \State $v \gets -\infty$
    \For{$a \in \mathit{moves}$}
      \State $v \gets \max\bigl(v,\ \textsc{Paranoid}(T(s,a),\,d{-}1,\,\alpha,\,\beta,\,\textsc{me})\bigr)$
      \State $\alpha \gets \max(\alpha, v)$
      \If{$\alpha \ge \beta$} \textbf{break} \EndIf
    \EndFor
  \Else \Comment{all other sides: minimising coalition}
    \State $v \gets +\infty$
    \For{$a \in \mathit{moves}$}
      \State $v \gets \min\bigl(v,\ \textsc{Paranoid}(T(s,a),\,d{-}1,\,\alpha,\,\beta,\,\textsc{me})\bigr)$
      \State $\beta \gets \min(\beta, v)$
      \If{$\alpha \ge \beta$} \textbf{break} \EndIf
    \EndFor
  \EndIf
  \State \Return $v$
\EndFunction
\end{algorithmic}
\end{algorithm}

\begin{table}[H]
  \centering
  \caption{AI parameters by game mode.}
  \label{tab:ai_params}
  \begin{tabular}{lccc}
    \toprule
    \textbf{Mode} & \textbf{Algorithm} & \textbf{Depth}
    & \textbf{Time/move} \\
    \midrule
    2-player & Minimax Alpha-Beta & 6 & $<0.3$\,s \\
    3-player & Paranoid Minimax   & 4 & $<0.5$\,s \\
    4-player & Paranoid Minimax   & 3 & $<0.5$\,s \\
    \bottomrule
  \end{tabular}
\end{table}

The AI runs in a background thread (\texttt{asyncio.run\_in\_executor}) so
that the interface does not freeze while it computes.

\subsection{Pruning effectiveness and move ordering}
\label{ssec:ordering}

The effectiveness of Alpha-Beta pruning depends strongly on the order in
which moves are visited. With branching factor $b$ and depth $d$, the
worst-case number of leaf nodes is $O(b^{d})$---no better than exhaustive
search---but if the best move is always visited first, this drops to
$O(b^{d/2})$, that is, it \emph{doubles} the feasible depth within the same
time budget~\citep{knuth1975alphabeta}. Move ordering is therefore the main
lever for pushing the depths in Table~\ref{tab:ai_params}.

Tactical Formation offers a cheap and natural ordering heuristic: because
each action \emph{foresees} the number of pieces it captures,
$|a|=|F_\delta|$ (Section~\ref{ssec:capture}), the bot visits moves in
\emph{decreasing} order of $|a|$. A move that captures many pieces is
usually strong and so is often the best move, helping pruning trigger
early. When $|a|$ is equal, the bot prefers a move that increases mobility
per Equation~\eqref{eq:eval}. To exploit this ordering fully, both engines
use \emph{make/unmake} of moves instead of copying the state at every node.
The 2-player engine additionally uses a \textbf{transposition table} keyed
by a Zobrist hash---storing the upper/lower/exact Alpha-Beta bounds and
prioritising the recorded best move---to avoid recomputing move sequences
that lead back to the same state; the Paranoid engine does not yet use a
transposition table because keying a state (together with the set of living
sides and the turn pointer) is more complex, left for a later version.
Iterative deepening is not yet implemented and is left for future work.

\subsection{Four design pitfalls and their fixes}
\label{ssec:pitfalls}

When testing the bot, we found four pitfalls that are \emph{not}
implementation bugs but stem from applying a general-purpose algorithm to a
specific ruleset: pieces are conserved (capture $=$ change of side,
Section~\ref{ssec:capture}) and the winning condition is essentially to
\emph{blockade} the opponent rather than to annihilate it. Each pitfall
reflects a wrong implicit assumption in that match. The first three pitfalls
(\#1--\#3) have concrete fixes that apply to both engines (2-player and
Paranoid); the fourth pitfall (\#4) is inherent to any depth-limited search
and is discussed separately at the end.

\paragraph{Pitfall \#1 --- Greedy capture.}
\emph{Cause.} The default move generator, when a move creates several
sandwich directions, keeps only the direction that flips the \emph{most}
pieces ($\arg\max_\delta |F_\delta|$). The bot never considers a smaller
capture, even when the large capture opens up a sandwich counterstrike.
Because pieces are conserved, a move that captures 5 pieces may be flipped
back by the opponent for $\ge 6$ pieces on the very next turn---exactly the
``bait'' that the sandwich rule allows.
\emph{Fix.} Replace the action set $\mathcal A(c,s)$ with the
\emph{expanded} set $\mathcal A^{+}$: each landing cell with $k\ge 2$
sandwich directions generates $k$ separate actions, one per direction. The
bot therefore considers smaller captures too; the search tree itself
(through the eval at deeper layers) decides how much capturing is
beneficial, instead of a fixed greedy rule.

\paragraph{Pitfall \#2 --- Wrong focus in the evaluation function.}
\emph{Cause.} The original evaluation function, $10(N_c-N_{\bar
c})+(M_c-M_{\bar c})$, weighted piece count ten times more heavily than
mobility. But because pieces are conserved, ``capturing everything'' is
nearly impossible; the real path to victory is to squeeze the opponent out
of moves. The bot thus optimised the wrong objective---hoarding pieces
instead of blockading. Moreover, the two terms have different units
(counting pieces versus counting moves), so the $10\!:\!1$ weight ratio
carries no clear meaning.
\emph{Fix.} Normalise both terms to $[-1,1]$ and then assign weights that
reflect a player's intuition: $0.7$ for capturing pieces, $0.3$ for
blocking the opponent's moves (Equations~\eqref{eq:eval}
and~\eqref{eq:evalpar}). A smaller weight for mobility but on the
\emph{same scale} makes it genuinely steer toward a blockade once the
piece-count advantage has saturated.

\paragraph{Pitfall \#3 --- Determinism and move repetition.}
\emph{Cause (two layers).} (i)~\emph{Determinism}: $\arg\max$ always
returns the first best move, so for the same position the bot always plays
identically---finding one winning line means winning identically forever.
(ii)~\emph{Oscillation}: because the eval is a continuous real number, two
moves almost never tie in absolute score, so the ``randomise on a tie''
strategy is useless. Worse, because pieces are conserved, the bot often
pushes a piece out and then pulls it straight back to the same position
($A\!\to\!B\!\to\!A$), forming a loop. Measured over 40 self-play games at
depth~2, each observed for at most its first 200 plies because without the
fix the bots can loop without ever finishing (5040 turns in total), \textbf{37.3\%} of positions were a position
that had \emph{already} appeared in the game---showing how common this
oscillation loop is without intervention. The key point: this pitfall
\emph{cannot} be fixed by intervening in the search---Alpha-Beta, move
ordering, or a transposition table change only the \emph{speed}, not the
\emph{move chosen}---so it must be fixed at the move-selection stage.
\emph{Fix (two mechanisms).}
\begin{itemize}
  \item \textbf{Anti-repetition penalty.} A game keeps a history
        $\mathcal H$ of the Zobrist hashes of \emph{every} position that
        has actually occurred. At the root, a move leading to a position
        $b'$ already in $\mathcal H$ is penalised in proportion to the
        number of repetitions:
        $\mathrm{eval}'(a)=\mathrm{eval}(a)-\lambda\cdot\#\{b'\in\mathcal
        H\}$ with $\lambda=0.5$. Because $\lambda$ is much larger than the
        eval gap between reasonable moves, the bot abandons the
        oscillation loop, yet it still accepts a repetition if that is the
        only move that does not lead to a loss.
  \item \textbf{Randomness within an $\varepsilon$ band.} Instead of an
        ``exact tie'', every move within a band $\varepsilon=0.02$ around
        the best move is treated as equivalent and chosen uniformly at
        random. This creates diversity between games without sacrificing
        quality---moves worse than the best by more than $\varepsilon$ are
        still rejected.
\end{itemize}
After the fix, repeating the same 40-game experiment gives a
position-repetition rate of \textbf{0.0\%} (0 repeated positions over 4069
turns)---the oscillation loop is eliminated entirely. Note that the history
$\mathcal H$ contains only positions that have \emph{already} occurred, so
it does not increase search cost exponentially; it is merely an $O(1)$ hash
lookup at the root.

\paragraph{Pitfall \#4 --- Horizon effect.}
A fourth pitfall in the development log---the \emph{horizon effect}, where a
danger is pushed beyond the search depth---is inherent to any depth-limited
search, not specific to this game. With a $5\times5$ board and depth 6, the
state space is small enough that this effect is negligible, so we leave it
as is rather than adding quiescence search; this is a natural direction for
future work if the board size is increased.
